% Options for packages loaded elsewhere
\PassOptionsToPackage{unicode}{hyperref}
\PassOptionsToPackage{hyphens}{url}
\documentclass[
  polish,
  11pt,
  a4paper,
]{article}
\usepackage{xcolor}
\usepackage[margin=2cm]{geometry}
\usepackage{amsmath,amssymb}
\setcounter{secnumdepth}{-\maxdimen} % remove section numbering
\usepackage{iftex}
\ifPDFTeX
  \usepackage[T1]{fontenc}
  \usepackage[utf8]{inputenc}
  \usepackage{textcomp} % provide euro and other symbols
\else % if luatex or xetex
  \usepackage{unicode-math} % this also loads fontspec
  \defaultfontfeatures{Scale=MatchLowercase}
  \defaultfontfeatures[\rmfamily]{Ligatures=TeX,Scale=1}
\fi
\usepackage{lmodern}
\ifPDFTeX\else
  % xetex/luatex font selection
\fi
% Use upquote if available, for straight quotes in verbatim environments
\IfFileExists{upquote.sty}{\usepackage{upquote}}{}
\IfFileExists{microtype.sty}{% use microtype if available
  \usepackage[]{microtype}
  \UseMicrotypeSet[protrusion]{basicmath} % disable protrusion for tt fonts
}{}
\makeatletter
\@ifundefined{KOMAClassName}{% if non-KOMA class
  \IfFileExists{parskip.sty}{%
    \usepackage{parskip}
  }{% else
    \setlength{\parindent}{0pt}
    \setlength{\parskip}{6pt plus 2pt minus 1pt}}
}{% if KOMA class
  \KOMAoptions{parskip=half}}
\makeatother
\usepackage{color}
\usepackage{fancyvrb}
\newcommand{\VerbBar}{|}
\newcommand{\VERB}{\Verb[commandchars=\\\{\}]}
\DefineVerbatimEnvironment{Highlighting}{Verbatim}{commandchars=\\\{\}}
% Add ',fontsize=\small' for more characters per line
\usepackage{framed}
\definecolor{shadecolor}{RGB}{248,248,248}
\newenvironment{Shaded}{\begin{snugshade}}{\end{snugshade}}
\newcommand{\AlertTok}[1]{\textcolor[rgb]{0.94,0.16,0.16}{#1}}
\newcommand{\AnnotationTok}[1]{\textcolor[rgb]{0.56,0.35,0.01}{\textbf{\textit{#1}}}}
\newcommand{\AttributeTok}[1]{\textcolor[rgb]{0.13,0.29,0.53}{#1}}
\newcommand{\BaseNTok}[1]{\textcolor[rgb]{0.00,0.00,0.81}{#1}}
\newcommand{\BuiltInTok}[1]{#1}
\newcommand{\CharTok}[1]{\textcolor[rgb]{0.31,0.60,0.02}{#1}}
\newcommand{\CommentTok}[1]{\textcolor[rgb]{0.56,0.35,0.01}{\textit{#1}}}
\newcommand{\CommentVarTok}[1]{\textcolor[rgb]{0.56,0.35,0.01}{\textbf{\textit{#1}}}}
\newcommand{\ConstantTok}[1]{\textcolor[rgb]{0.56,0.35,0.01}{#1}}
\newcommand{\ControlFlowTok}[1]{\textcolor[rgb]{0.13,0.29,0.53}{\textbf{#1}}}
\newcommand{\DataTypeTok}[1]{\textcolor[rgb]{0.13,0.29,0.53}{#1}}
\newcommand{\DecValTok}[1]{\textcolor[rgb]{0.00,0.00,0.81}{#1}}
\newcommand{\DocumentationTok}[1]{\textcolor[rgb]{0.56,0.35,0.01}{\textbf{\textit{#1}}}}
\newcommand{\ErrorTok}[1]{\textcolor[rgb]{0.64,0.00,0.00}{\textbf{#1}}}
\newcommand{\ExtensionTok}[1]{#1}
\newcommand{\FloatTok}[1]{\textcolor[rgb]{0.00,0.00,0.81}{#1}}
\newcommand{\FunctionTok}[1]{\textcolor[rgb]{0.13,0.29,0.53}{\textbf{#1}}}
\newcommand{\ImportTok}[1]{#1}
\newcommand{\InformationTok}[1]{\textcolor[rgb]{0.56,0.35,0.01}{\textbf{\textit{#1}}}}
\newcommand{\KeywordTok}[1]{\textcolor[rgb]{0.13,0.29,0.53}{\textbf{#1}}}
\newcommand{\NormalTok}[1]{#1}
\newcommand{\OperatorTok}[1]{\textcolor[rgb]{0.81,0.36,0.00}{\textbf{#1}}}
\newcommand{\OtherTok}[1]{\textcolor[rgb]{0.56,0.35,0.01}{#1}}
\newcommand{\PreprocessorTok}[1]{\textcolor[rgb]{0.56,0.35,0.01}{\textit{#1}}}
\newcommand{\RegionMarkerTok}[1]{#1}
\newcommand{\SpecialCharTok}[1]{\textcolor[rgb]{0.81,0.36,0.00}{\textbf{#1}}}
\newcommand{\SpecialStringTok}[1]{\textcolor[rgb]{0.31,0.60,0.02}{#1}}
\newcommand{\StringTok}[1]{\textcolor[rgb]{0.31,0.60,0.02}{#1}}
\newcommand{\VariableTok}[1]{\textcolor[rgb]{0.00,0.00,0.00}{#1}}
\newcommand{\VerbatimStringTok}[1]{\textcolor[rgb]{0.31,0.60,0.02}{#1}}
\newcommand{\WarningTok}[1]{\textcolor[rgb]{0.56,0.35,0.01}{\textbf{\textit{#1}}}}
\usepackage{longtable,booktabs,array}
\usepackage{caption}
\captionsetup[table]{skip=6pt}
\newcounter{none} % for unnumbered tables
\usepackage{calc} % for calculating minipage widths
% Correct order of tables after \paragraph or \subparagraph
\usepackage{etoolbox}
\makeatletter
\patchcmd\longtable{\par}{\if@noskipsec\mbox{}\fi\par}{}{}
\makeatother
% Allow footnotes in longtable head/foo
\IfFileExists{footnotehyper.sty}{\usepackage{footnotehyper}}{\usepackage{footnote}}
\makesavenoteenv{longtable}
\usepackage{graphicx}
\makeatletter
\newsavebox\pandoc@box
\newcommand*\pandocbounded[1]{% scales image to fit in text height/width
  \sbox\pandoc@box{#1}%
  \Gscale@div\@tempa{\textheight}{\dimexpr\ht\pandoc@box+\dp\pandoc@box\relax}%
  \Gscale@div\@tempb{\linewidth}{\wd\pandoc@box}%
  \ifdim\@tempb\p@<\@tempa\p@\let\@tempa\@tempb\fi% select the smaller of both
  \ifdim\@tempa\p@<\p@\scalebox{\@tempa}{\usebox\pandoc@box}%
  \else\usebox{\pandoc@box}%
  \fi%
}
% Set default figure placement to htbp
\def\fps@figure{htbp}
\makeatother
\ifLuaTeX
\usepackage[bidi=basic,shorthands=off]{babel}
\else
\usepackage[bidi=default,shorthands=off]{babel}
\fi
\ifLuaTeX
  \usepackage{selnolig} % disable illegal ligatures
\fi
\setlength{\emergencystretch}{3em} % prevent overfull lines
\providecommand{\tightlist}{%
  \setlength{\itemsep}{0pt}\setlength{\parskip}{0pt}}
\usepackage{fontspec}
\defaultfontfeatures{Ligatures=TeX}
\usepackage{amsmath,amssymb}
\usepackage{booktabs}
\usepackage{array}
\usepackage{geometry}
\usepackage{enumitem}
\usepackage{xcolor}
\usepackage{tcolorbox}
\usepackage{fancyhdr}
\usepackage{titlesec}
\usepackage{multicol}

\geometry{margin=1.8cm, top=2.5cm, bottom=2cm}
\pagestyle{fancy}
\fancyhf{}
\rhead{\small Projektowanie badań — 2026/27}
\lhead{\small Zarządzanie informacją | ISI UJ}
\cfoot{\thepage}

\titleformat{\section}{\large\bfseries\color{blue!70!black}}{\thesection.}{0.5em}{}
\titleformat{\subsection}{\normalsize\bfseries}{\thesubsection.}{0.5em}{}
\titleformat{\subsubsection}{\small\bfseries\color{gray!70!black}}{\thesubsubsection.}{0.5em}{}
\titlespacing*{\section}{0pt}{1.5ex}{1ex}
\titlespacing*{\subsection}{0pt}{1ex}{0.5ex}

\setlist[itemize]{nosep, leftmargin=1.5em}
\setlist[enumerate]{nosep, leftmargin=1.5em}

\tcbset{boxrule=0.5pt, left=4pt, right=4pt, top=4pt, bottom=4pt, fonttitle=\bfseries\small}

\newtcolorbox{definicja}[1][]{colback=blue!5, colframe=blue!50!black, title={#1}}
\newtcolorbox{uwagabox}[1][]{colback=orange!10, colframe=orange!70!black, title={#1}}
\newtcolorbox{wzorbox}[1][]{colback=gray!5, colframe=gray!50!black, title={#1}}
\newtcolorbox{przykladbox}[1][]{colback=purple!5, colframe=purple!50!black, title={#1}}

\usepackage{bookmark}
\IfFileExists{xurl.sty}{\usepackage{xurl}}{} % add URL line breaks if available
\urlstyle{same}
% fallback for those not using the hyperref driver hyperxmp:
\makeatletter
\@ifundefined{xmpquote}{\newcommand{\xmpquote}[1]{#1}}{}
\makeatother
\hypersetup{
  pdftitle={Handout C02 --- Import i porządkowanie danych ankietowych},
  pdfauthor={Marek Deja},
  pdflang={pl-PL},
  hidelinks,
  pdfcreator={LaTeX via pandoc}}

\title{Handout C02 --- Import i porządkowanie danych ankietowych}
\usepackage{etoolbox}
\makeatletter
\providecommand{\subtitle}[1]{% add subtitle to \maketitle
  \apptocmd{\@title}{\par {\large #1 \par}}{}{}
}
\makeatother
\subtitle{Zarządzanie informacją --- część ilościowa --- 2026/27}
\author{Marek Deja}
\date{2026/27}

\begin{document}
\maketitle

\section{C02 --- plan 90 minut}\label{c02-plan-90-minut}

10 min login/odtworzenie → 20 min import i typy → 25 min wspólne
czyszczenie → 25 min Z02 → 10 min wysyłka.

\section{Kontrola po imporcie}\label{kontrola-po-imporcie}

CSV kursu: przecinek jako separator, kropka dziesiętna, UTF-8. Odczytamy
tabelę i jej strukturę. Sprawdź N, kolumny i etykiety; identyfikator nie
jest miarą ilościową.

\begin{Shaded}
\begin{Highlighting}[]
\NormalTok{surowe }\OtherTok{\textless{}{-}}\NormalTok{ badaniaZI}\SpecialCharTok{::}\FunctionTok{dane\_przykladowe}\NormalTok{()}
\FunctionTok{dim}\NormalTok{(surowe)}
\end{Highlighting}
\end{Shaded}

\begin{verbatim}
## [1] 150  15
\end{verbatim}

\begin{Shaded}
\begin{Highlighting}[]
\FunctionTok{str}\NormalTok{(surowe[, }\FunctionTok{c}\NormalTok{(}\StringTok{"id\_odpowiedzi"}\NormalTok{, }\StringTok{"grupa"}\NormalTok{, }\StringTok{"czas\_wyszukiwania"}\NormalTok{)])}
\end{Highlighting}
\end{Shaded}

\begin{verbatim}
## 'data.frame':    150 obs. of  3 variables:
##  $ id_odpowiedzi    : chr  "o001" "o002" "o003" "o004" ...
##  $ grupa            : chr  "nowi" "nowi" "doświadczeni" "doświadczeni" ...
##  $ czas_wyszukiwania: num  12.1 10.4 5.8 16.1 28 6.3 8.1 24.4 18.4 18.7 ...
\end{verbatim}

{\def\LTcaptype{none} % do not increment counter
\begin{longtable}[]{@{}ll@{}}
\toprule\noalign{}
Zmienna & Skala i reguła
\midrule\noalign{}
\endhead
\bottomrule\noalign{}
\endlastfoo
grupa & nominalna, etykiety/factor
czas\_wyszukiwania & ilościowa, minuty, zakres 0--120
pozycja\_1--6 & porządkowa 1--5
pozycja\_2 = 99 & brak odpowiedzi, nie wysoka ocena
kanal\_* = 0 & kanał niezaznaczony, nie brak
NA & wartość nieznana
\end{longtable}
}

\section{Czyszczenie --- kopia danych}\label{czyszczenie-kopia-danych}

Usuń tylko identyczne kopie całych wierszy. Rekoduj opisany w słowniku
kod i czas spoza zakresu w konkretnych komórkach. Odczytaj N po
czyszczeniu; przy sprzecznym ID wyjaśnij konflikt.

\begin{Shaded}
\begin{Highlighting}[]
\NormalTok{dane }\OtherTok{\textless{}{-}}\NormalTok{ surowe[}\SpecialCharTok{!}\FunctionTok{duplicated}\NormalTok{(surowe), ]}
\NormalTok{k }\OtherTok{\textless{}{-}} \SpecialCharTok{!}\FunctionTok{is.na}\NormalTok{(dane}\SpecialCharTok{$}\NormalTok{pozycja\_2) }\SpecialCharTok{\&}\NormalTok{ dane}\SpecialCharTok{$}\NormalTok{pozycja\_2 }\SpecialCharTok{==} \DecValTok{99}
\NormalTok{dane}\SpecialCharTok{$}\NormalTok{pozycja\_2[k] }\OtherTok{\textless{}{-}} \ConstantTok{NA}
\NormalTok{k }\OtherTok{\textless{}{-}} \SpecialCharTok{!}\FunctionTok{is.na}\NormalTok{(dane}\SpecialCharTok{$}\NormalTok{czas\_wyszukiwania) }\SpecialCharTok{\&}\NormalTok{ (dane}\SpecialCharTok{$}\NormalTok{czas\_wyszukiwania }\SpecialCharTok{\textless{}} \DecValTok{0} \SpecialCharTok{|}\NormalTok{ dane}\SpecialCharTok{$}\NormalTok{czas\_wyszukiwania }\SpecialCharTok{\textgreater{}} \DecValTok{120}\NormalTok{)}
\NormalTok{dane}\SpecialCharTok{$}\NormalTok{czas\_wyszukiwania[k] }\OtherTok{\textless{}{-}} \ConstantTok{NA}
\FunctionTok{nrow}\NormalTok{(dane)}
\end{Highlighting}
\end{Shaded}

\begin{verbatim}
## [1] 148
\end{verbatim}

Zapisz liczby zmian \textbf{przed} rekodowaniem. Dziennik rozróżnia
wiersze i komórki. Nie zmieniaj \texttt{dane/surowe.csv}; eksportuj nowy
CSV bez numerów wierszy.

\section{Z02 --- dowody i punkty}\label{z02-dowody-i-punkty}

Własny import, polskie znaki, N i klasy --- 2 pkt. Typy, 99/NA i
identyczne duplikaty --- 4 pkt. Dziennik reguł/liczb/N i rozróżnienie
brak/zero --- 2 pkt. \texttt{wyczyszczone.csv}, \texttt{dziennik.csv} i
działający od początku skrypt --- 2 pkt.

Pliki: \texttt{zadania/z02/analiza.R}, \texttt{odpowiedzi.md},
\texttt{wyniki/wyczyszczone.csv}, \texttt{wyniki/dziennik.csv}. W
\texttt{.md} wyjaśnij reguły i dodaj względne linki do wyników. Po
zapisie uruchom \texttt{sprawdz\_zadanie("Z02")},
\texttt{oddaj\_zadanie("Z02")}, \texttt{status\_oddania("Z02")}; zakończ
przez \texttt{wyloguj\_github()} i wylogowanie przeglądarki.

\section{Notatki}\label{notatki}

N przed: \ldots\ldots\ldots. N po: \ldots\ldots\ldots. Kopie:
\ldots\ldots\ldots. Kod 99: \ldots\ldots\ldots. Czasy:
\ldots\ldots\ldots.

Dlaczego nie usuwam osoby za jeden brak?
\ldots\ldots\ldots\ldots\ldots\ldots\ldots\ldots\ldots\ldots\ldots\ldots\ldots\ldots\ldots\ldots\ldots\ldots..

Dlaczego klasa liczbowa nie wystarcza do uznania skali za ilościową?
\ldots\ldots\ldots\ldots\ldots\ldots\ldots\ldots\ldots\ldots..

\end{document}
